\section{Literature Review}

The reading for this project covered three strands: what is known about
augmenting instance segmentation, what is known about using generative models to
fill a long tail, and what the floor-plan and foundation-model literature offers
as backbones and datasets. Each paper below is summarised with its venue,
availability of code and data, method, main finding, and a comment on what it
implied for the design of this project.

\paragraph{Ghiasi et al., ``Simple Copy-Paste is a Strong Data Augmentation
Method for Instance Segmentation'' (CVPR 2021)
\cite{ghiasi2021copypaste}.} Code is released inside the TensorFlow TPU model
repository; the paper uses public COCO and LVIS data. The method is deliberately
minimal: instance masks are cut from one training image and pasted at random
scale and position into another, with no blending model and no attempt to make
the composite photometrically consistent. On COCO the authors report $+1.0$ box
AP over a strong baseline, and on LVIS the gain concentrates in the rare
categories ($+6.1$ mask AP for rare classes), while remaining complementary to a
long-tail loss. The finding that a crude, generator-free pixel transplant
outperforms more sophisticated composition is the single most influential result
for this project: it motivated including direct copy-paste as an explicit control
arm (e8) rather than only as an ablation, and that arm turned out to be the only
augmentation with a consistent positive sign in our own measurements.

\paragraph{Xie et al., ``MosaicFusion: Diffusion Models as Data Augmenters for
Large Vocabulary Instance Segmentation'' (IJCV 2024; arXiv:2309.13042)
\cite{xie2024mosaicfusion}.} Code and generated data are public at
\url{https://github.com/Jiahao000/MosaicFusion}. The method is training-free: an
off-the-shelf text-to-image diffusion model generates a canvas divided into
regions, each conditioned on a different category prompt, and the instance masks
are read out of the cross-attention maps rather than annotated. On LVIS the
authors report consistent gains for rare and common categories with several
detector architectures. The paper establishes the strongest form of the
hypothesis this project set out to test --- that a general-purpose diffusion
model can manufacture usable training instances for categories the detector has
barely seen. The critical difference in our setting is domain: MosaicFusion
operates on natural images, where the pretrained diffusion prior is a good model
of the data distribution, whereas a floor plan is line art the base model was
never trained to draw. Our results suggest that difference is decisive.

\paragraph{Shin et al., ``Fill-Up: Balancing Long-Tailed Data with Generative
Models'' (arXiv:2306.07200, 2023) \cite{shin2023fillup}.} Code and a project
page are available at \url{https://alex4727.github.io/Fill-Up/}. Rather than
prompting a frozen model with class names, the authors learn a textual-inversion
token per tail category from the few real examples available, then generate
synthetic images from that token until the class distribution is balanced. They
report that the inverted tokens align the synthetic images with the real domain
far better than plain prompts, and that filling the tail this way reaches
state-of-the-art results on standard long-tailed classification benchmarks when
combined with existing long-tail techniques. The paper is the direct argument for
the LoRA branch of this project: if plain prompting fails to place a class in the
right visual domain, then adapting the generator to the domain should fix it.
Our LoRA arms tested that argument on floor plans and did not reproduce the
benefit, which is worth stating precisely because the reasoning is sound and the
prior evidence is positive.

\paragraph{Wang et al., ``Domain Gap Embeddings for Generative Dataset
Augmentation'' (CVPR 2024) \cite{wang2024doge}.} Code is at
\url{https://github.com/humansensinglab/DoGE}. The method extracts the
difference between a source and a desired data distribution as a latent
embedding, then uses that embedding to steer a frozen generative model, avoiding
per-dataset fine-tuning entirely. It is evaluated under subpopulation shift and
three few-shot domain-adaptation scenarios, improving over both classical
augmentation and fine-tuned generative augmentation. The relevant lesson is
diagnostic: DoGE is built around the observation that generative augmentation
underperforms precisely when the target distribution is far from the generator's
own, which is a good description of the failure mode we observed. It also
suggests the most plausible next step for this project, since the method needs no
adapter training and would be applicable to a 30-example tail.

\paragraph{Huang et al., ``Gen-n-Val: Agentic Image Data Generation and
Validation'' (arXiv:2506.04676, 2025; CVPR 2026 Findings)
\cite{huang2025gennval}.} Code is linked from the arXiv page. The method wraps
layer-diffusion generation in an agentic loop in which a large language model
writes and revises the prompts and a vision-language model validates each
generated instance mask before it is admitted to the training set. The reported
headline is that the proportion of invalid synthetic samples drops from roughly
50\% to 7\%, with corresponding gains on rare-class instance segmentation and
open-vocabulary detection. The paper is the clearest statement in the recent
literature that the bottleneck in generative augmentation is acceptance, not
generation. Our own pipeline enforces acceptance through deterministic raster
checks --- exact invariance outside the edited region, retention of donor
strokes, pixel-statistics thresholds --- which certify structural safety but
cannot certify that the generated interior \emph{means} the requested room type.
Gen-n-Val's semantic validator is the component we lacked, and its absence is a
candidate explanation for our null result.

\paragraph{Song et al., ``TMI: Text-to-Image Meets Image-to-Image for
Complementary Data Synthesis to Boost Long-Tailed Instance Segmentation''
(arXiv:2607.08201, 2026; ECCV 2026) \cite{song2026tmi}.} Code is linked from the
arXiv page. The method combines two synthesis routes that fail differently:
text-to-image generation, which produces diverse instances but with weak
correspondence to a usable mask, and image-to-image translation of existing
annotated instances, which preserves the mask but produces little diversity. The
authors show the two are complementary on long-tailed instance-segmentation
benchmarks and that using either alone leaves most of the benefit unrealised.
This maps almost exactly onto the axis explored in our Strategy A/B (prompt-driven,
diverse, semantically unanchored) versus Strategy C (exemplar-driven,
mask-faithful, low-diversity) arms, and it predicts that neither family should
work well alone --- which is what we measured, although our experiment does not
include the combination the paper recommends.

\paragraph{Wang et al., ``Seesaw Loss for Long-Tailed Instance Segmentation''
(CVPR 2021) \cite{wang2021seesaw}.} Code is in MMDetection. Instead of adding
data, the method rebalances the classification loss with two multiplicative
factors: a mitigation factor that suppresses gradients a frequent class exerts on
a rare one, and a compensation factor that increases the penalty for
misclassifying a rare instance. It reaches a substantial rare-AP improvement on
LVIS with no change to the data pipeline. The paper is the strongest reminder
that data synthesis is not the only lever, and that a loss-side treatment is far
cheaper than a diffusion pipeline. Given our finding that the classification
stage, not the localisation stage, is the failure point, a Seesaw-style loss is
arguably the intervention this project should have tried first, and it is the
main untested alternative for future work.

\paragraph{Gupta et al., ``LVIS: A Dataset for Large Vocabulary Instance
Segmentation'' (CVPR 2019) \cite{gupta2019lvis}.} Data and evaluation code are at
\url{https://www.lvisdataset.org/}. LVIS introduced the frequent/common/rare
split protocol and the federated annotation scheme that every long-tail
instance-segmentation paper above evaluates against. It matters here mainly as a
scale reference: LVIS defines ``rare'' as fewer than 10 training images per
category across a set of 100k images, whereas our rare categories have 3--8
training instances in a set of 35 images. The methods surveyed above are
therefore validated in a regime that has two to three orders of magnitude more
data than ours, which is a caution the results section returns to.

\paragraph{Kalervo et al., ``CubiCasa5K: A Dataset and an Improved Multi-Task
Model for Floorplan Image Analysis'' (SCIA 2019)
\cite{kalervo2019cubicasa5k}.} Data and code are at
\url{https://github.com/CubiCasa/CubiCasa5k}. The paper contributes 5000
annotated Finnish residential floor plans over more than 80 object categories and
a multi-task network that predicts room segmentation, icon segmentation and wall
heatmaps from a shared encoder. It remains the standard reference dataset for
raster floor-plan parsing. Its relevance here is as a contrast: CubiCasa5K's room
classes are residential and visually distinctive, its scale is two orders of
magnitude above ours, and its plans follow a consistent drawing convention. Our
dataset shares none of those properties, which is why transferring published
architectures was never expected to be sufficient on its own.

\paragraph{Kratochvila et al., ``Recognizing and Reconstructing a Multi-Unit
Floor Plan'' (arXiv:2408.01526, 2024) \cite{kratochvila2024multiunit}.} The
paper addresses raster multi-unit plans directly, comparing U-Net variants
(MDA-UNet, MACU-Net) for pixel-wise segmentation of floor-plan structure and
reporting an F1 around 0.86 on their target elements, followed by a
vectorisation step. It is the closest recent work to our input domain, and it is
informative for two reasons. First, it confirms that structural elements ---
walls and room boundaries --- are the tractable part of the problem, which
matches our class-agnostic AP of 0.59--0.71. Second, it does not attempt
fine-grained functional room typing over a long tail, which is the part we found
to be the hard one, so it offers no evidence that the harder task is reachable at
our scale.

\paragraph{Ganon et al., ``WAFFLE: Multimodal Floorplan Understanding in the
Wild'' (arXiv:2412.00955; WACV 2025) \cite{ganon2025waffle}.} Data and code are
at \url{https://tau-vailab.github.io/WAFFLE/}. The authors assemble roughly 20k
floor plans scraped from the web, spanning building types, drawing conventions
and languages, and use large language and vision-language models to extract
semantic labels from the heterogeneous text on the drawings. They demonstrate
semantic floor-plan tasks that single-convention datasets cannot support. The
paper is the closest published example of foundation models being applied to
architectural floor plans in a few-shot manner, and it points at a resource this
project did not use: our plans carry printed German room labels
(``B\"uro'', ``Technik'') that a text-reading model could exploit directly. Given
that our vision-only classifier reached zero rare recall, reading the drawing's
own text is probably a stronger signal than anything a diffusion model can add.

\paragraph{Carion et al., ``SAM 3: Segment Anything with Concepts''
(arXiv:2511.16719, 2025) \cite{carion2025sam3}.} Code and weights are at
\url{https://github.com/facebookresearch/sam3}, gated on Hugging Face. SAM 3
extends promptable segmentation to \emph{concept} prompts: a short noun phrase
retrieves and segments all matching instances in an image, backed by a detector
with a presence head that decouples ``is this concept present'' from ``where is
it''. The model and its official trainer are the reason Goal 3 was feasible at
all --- fine-tuning it on 35 images with 19 German category names as text
concepts required a configuration change rather than a new architecture. In our
experiments it delivered the single largest effect measured in the project, which
reframed the augmentation question as a second-order one.

\paragraph{Oquab et al., ``DINOv2: Learning Robust Visual Features without
Supervision'' (TMLR 2024; arXiv:2304.07193) \cite{oquab2024dinov2}.} Code and
weights are at \url{https://github.com/facebookresearch/dinov2}. The paper
trains ViT backbones by self-distillation on a curated 142M-image corpus and
shows that the resulting frozen features support linear-probe classification,
depth estimation and segmentation at a level previously requiring supervised
fine-tuning. The practical claim exploited in Goal 4 is that a linear head on
frozen features is competitive when labelled data is scarce, which is exactly our
regime. The pretraining corpus is natural images, however, and our result ---
0.45 top-1 over 19 classes and zero recall on all six rare classes --- suggests
the frozen features do not separate functional room types in line-art drawings.

\paragraph{Zhuo et al., ``UINO-FSS: Unifying Representation Learning and
Few-Shot Segmentation via Hierarchical Distillation and Mamba-HyperCorrelation''
(arXiv:2504.15669, 2025) \cite{zhuo2025uinofss}.} The authors observe that
early-stage DINOv2 features are distributionally consistent with SAM's output
embeddings, and exploit this to distil a single backbone that serves both
representation learning and few-shot segmentation, reporting mIoU gains of
3.8--4.1 points on standard few-shot benchmarks. The paper is the strongest
recent evidence that a frozen DINOv2 encoder is a viable few-shot segmentation
substrate, and it is the reason we tested a frozen encoder rather than
fine-tuning one. Its benchmarks (PASCAL-$5^i$, COCO-$20^i$) use natural-image
classes with clear appearance signatures, which the negative outcome of our
Goal 4 experiment suggests is a load-bearing assumption.

\paragraph{Rombach et al., ``High-Resolution Image Synthesis with Latent
Diffusion Models'' (CVPR 2022) \cite{rombach2022ldm}.} Code is at
\url{https://github.com/CompVis/latent-diffusion} and the derived Stable
Diffusion 1.5 inpainting checkpoint is on Hugging Face. Running the diffusion
process in the latent space of a pretrained autoencoder rather than in pixel
space makes high-resolution synthesis and inpainting tractable on a single GPU,
and the inpainting variant conditions the denoiser on a mask and the masked
image through a nine-channel input. This is the generator used by every
diffusion arm in this project. The paper's own inpainting demonstrations are
photographic, and the qualitative failures we observed --- perspective furniture
and photographic texture rendered inside a 2D line drawing, visible in
Figures~\ref{fig:qastrategya} and \ref{fig:loraval} --- are consistent with a
prior that has essentially no floor-plan mass.

\paragraph{Hu et al., ``LoRA: Low-Rank Adaptation of Large Language Models''
(ICLR 2022) \cite{hu2022lora}.} Code is at
\url{https://github.com/microsoft/LoRA}. The method freezes the pretrained
weights and learns a low-rank update $BA$ for selected projection matrices,
reducing trainable parameters by orders of magnitude while matching full
fine-tuning quality, and adding no inference latency once merged. It has become
the standard way to adapt a diffusion U-Net to a new visual domain on a small
dataset, which is precisely what the LoRA branch of this project attempted with
rank 8 on the attention projections. The technique itself worked as advertised ---
training converged and the adapter demonstrably changed the generated
output --- but Section~\ref{sec:results} shows the adapted generator did not
produce more useful training data, and on the class-agnostic track it produced
the single worst arm in the grid.
